\documentclass[11pt,a4paper]{article}
\usepackage[margin=1in]{geometry}
\usepackage{amsmath}
\usepackage{hyperref}
\usepackage{setspace}

\setlength{\parskip}{0.5em}
\setlength{\parindent}{0pt}
\onehalfspacing

\title{Computer Animation HW3 Report\\Inverse Kinematics (Jacobian and CCD)}
\author{Student ID: 112550052}
\date{\today}

\begin{document}
\maketitle

\section{Question: Local and Global Coordinates (5\%)}

\textbf{Global (world) coordinates} describe geometry in a single fixed frame shared by the whole scene.
In this assignment, each bone stores world-space quantities after forward kinematics, for example
\texttt{start\_position}, \texttt{end\_position}, and the accumulated orientation \texttt{rotation}.
Targets for IK are also given in this same world frame.

\textbf{Local coordinates} describe vectors and rotations relative to a moving frame attached to a bone
(or to the joint frame defined by the ASF hierarchy and the \texttt{rot\_parent\_current} offset).
The AMC posture stores per-bone Euler angles (\texttt{Posture::bone\_rotations}) that act as joint variables
in the bone's local rotation convention (ZYX composition via \texttt{util::rotateDegreeZYX} in the provided FK).
The bone direction \texttt{dir} is expressed in the bone's local frame before it is transformed by the
current world rotation when computing the bone endpoint.

\textbf{Relationship.}
Forward kinematics maps local joint parameters to global poses by walking the hierarchy and concatenating
transforms. Inverse kinematics typically computes a desired correction in world space (for example aligning
the vector from a joint to the effector toward the target), then maps that correction back into changes to
local Euler parameters, applies joint limits, and runs FK again so the next iteration sees consistent globals.

\section{Implementation (10\%)}

\subsection{\texttt{IKConverged}}
For each IK chain index $i$, the end effector position is read through \texttt{jointChains[$i$][0]} (a pointer
into the live bone world position). Convergence holds when every chain satisfies
$\lVert \mathbf{p}_{\mathrm{ee},i}-\mathbf{t}_i\rVert < \varepsilon$ with $\varepsilon = 10^{-3}$.

\subsection{\texttt{clampBoneRotation}}
Rotation limits from the ASF file (\texttt{rxmin/rxmax}, etc.) are applied only to enabled degrees of freedom
(\texttt{dofrx}, \texttt{dofry}, \texttt{dofrz}). Disabled axes are left unchanged.

\subsection{\texttt{pseudoInverseLinearSolver}}
The Jacobian linearization is solved with a damped thin SVD pseudoinverse, minimizing
$\lVert J\Delta\boldsymbol{\theta}-\mathbf{e}\rVert$ while down-weighting near-singular directions.
The damping factor is tunable at runtime (bonus; default $10^{-6}$).

\subsection{\texttt{JacobianChainSolver}}
Each iteration runs FK, forms the geometric Jacobian for positional rows using instantaneous world axes for
the ZYX Euler parameterization (matching the course slides), scales the position error by the (runtime-tunable)
Jacobian step, solves for $\Delta\boldsymbol{\theta}$, converts radians to degrees, applies updates, clamps,
applies the root translation fix for chain index $0$, and runs FK again.

\subsection{\texttt{CCDChainSolver}}
Each outer iteration begins with FK and an early exit if the effector is close to the target.
CCD sweeps joints from the end of the chain toward the base. At each joint pivot, normalized vectors to the
effector and to the target define a shortest quaternion rotation in world space; its angle is scaled by the
(runtime-tunable) CCD step. The corresponding spatial angular increment is mapped to local Euler rates using the
same instantaneous axis columns as the Jacobian orientation part, via a small damped least-squares solve
(damping tunable at runtime for the bonus), then updates are clamped and FK refreshes the effector pointer
\texttt{jointChains[0][0]} for the next joint.

\subsection{\texttt{inverseIKSolver}}
FK is run once before iterating chains so shared postures begin from a consistent world state. Each chain
solver updates the same posture; if the convergence test fails, the posture is restored from a copy and FK is
run again.

\section{Discussion (5\%)}

\subsection{Step size}
Jacobian and CCD step sizes trade stability against speed.
Smaller values reduce overshoot and oscillation near targets and near joint limits, but require more iterations.
Larger values can track moving targets faster but may violate limits more aggressively before clamping, and
can diverge when the linearization or single-joint alignment is inaccurate.
The bonus UI exposes these as sliders so the trade-off can be explored interactively.

\subsection{Jacobian versus CCD}
The Jacobian method updates all joints in the chain together using a global linearization of the end effector
motion, which tends to behave smoothly for small residuals and is closer to the classical inverse-Jacobian
pipeline from lecture, but it can be sensitive to singularities and coupling between Euler axes.
Increasing SVD damping in the Jacobian solve often smooths motion near singular configurations at the cost of
slower convergence.
CCD applies localized joint-by-joint corrections and is often robust and easy to implement, but the sweep
ordering matters and convergence can be slower or exhibit characteristic ``CCD'' motion under large errors.

\section{Bonus (optional)}

I added \textbf{runtime IK parameter tuning} in the Control panel under \textbf{``Bonus: IK tuning''}:
\begin{itemize}
    \item \textbf{Jacobian step} and \textbf{CCD step} scale how aggressively each solver moves per iteration.
    \item \textbf{Jacobian SVD damping} softens the Moore--Penrose solve for $J\Delta\theta \approx e$ near singularities.
    \item \textbf{CCD axis-fit damping} regularizes the $3\times k$ least-squares mapping from world angular increment
          to Euler rates at each CCD joint.
    \item \textbf{Reset IK parameters} restores course defaults ($0.3$, $0.3$, $10^{-6}$, $10^{-4}$).
\end{itemize}
API: \texttt{kinematics.h} declares getters, setters, and \texttt{resetIKParametersToDefault()}; defaults live in
\texttt{kinematics.cpp} (\texttt{internal} namespace). The ImGui controls are in \texttt{InverseKinematics/main.cpp}.

\section{Miscellaneous}
Modified files for the bonus: \texttt{include/simulation/kinematics.h}, \texttt{src/simulation/kinematics.cpp},
\texttt{InverseKinematics/main.cpp}. Core IK remains in \texttt{kinematics.cpp}.

\end{document}
